\input{ieee-preamble}
\hypersetup{
  pdftitle={Underactuated Humanoid Locomotion},
  pdfauthor={Santiago Restrepo},
  pdfsubject={Simulation-first approach to learned humanoid walking and underactuation},
  pdfkeywords={humanoid, locomotion, underactuation, MuJoCo, C++, imitation, PPO}
}

\title{Underactuated Humanoid Locomotion}
\author{\IEEEauthorblockN{Santiago Restrepo}
\IEEEauthorblockA{\href{mailto:contact@waajacu.com}{contact@waajacu.com}\quad
\url{https://waajacu.com/}\\[3pt]
\textit{Learning-based control method for bipedal walking}}}
\date{}

\begin{document}
\maketitle

\begin{abstract}
This brief describes a simulation-first approach to humanoid walking and
underactuated motion. A published walking controller provides a baseline for
local policy development through imitation and reinforcement learning.
Physical diagnostics and visual inspection guide bounded experiments. The
immediate task is recognizable walking; deliberately passive joints follow.
\end{abstract}

\begin{IEEEkeywords}
Humanoid locomotion, underactuation, simulation, reinforcement learning.
\end{IEEEkeywords}

\balance
\section{Model and Execute the Dynamics}
Use a free base, actuated legs, and fixed upper-body joints. For configuration
$q$ and generalized velocity $v$, the dynamics are~\cite{tedrakeContact}
\begin{equation}
\begin{aligned}
M(q)\dot v+h(q,v)&=B\tau+J_c^{\mathsf T}\lambda,\\
\dot q&=N(q)v.
\end{aligned}
\end{equation}
$M$ is inertia, $h$ includes gravity and velocity terms, $B$ maps motor
torques $\tau$, $J_c^{\mathsf T}$ maps contact forces $\lambda$, and $N$
handles orientation kinematics. The base is unactuated; contact changes control
authority. Fixed flat-foot support can be fully actuated in reduced coordinates.

Native C++ connects MuJoCo contact and friction dynamics to LibTorch learning
and inference. Physics runs on the CPU; neural computation uses CUDA. A browser
shows live simulation frames and motion/contact diagnostics.

\section{Establish the Walking Baseline}
Replay the published Unitree controller~\cite{unitreeGym} with its documented
observation and action interface. Preserve it as a reference alongside a
zero-torque control. The local PPO actor receives command $c$ and causal history
$H$ of angular velocity, projected gravity, joint states, and prior actions:
\begin{equation}
\begin{aligned}
a_t&\sim\pi_\theta(\,\cdot\mid H_t,c_t),\\
q_{\mathrm{des}}&=q_{\mathrm{nom}}+D\tanh(a_t),\\
\tau&=\operatorname{sat}[K_p(q_{\mathrm{des}}-q_a)-K_dv_a].
\end{aligned}
\end{equation}
$\theta$ denotes policy parameters and $D$ bounds nominal-pose offsets.
Joint states $q_a$, $v_a$ drive faster proportional-derivative (PD) feedback;
$\operatorname{sat}$ enforces motor limits. Preserve preprocessing and action
transformations when exporting and evaluating each policy.

An asymmetric critic receives additional simulator state during training.
The actor uses observable history without privileged critic signals.

\section{Develop and Adapt Local Policies}
Imitation trains a student from the reference controller's action labels.
Independent reinforcement learning starts from random weights without teacher
labels. Proximal policy optimization (PPO)~\cite{schulman2017} collects parallel
rollouts, estimates advantages, and applies clipped updates to optimize
\begin{equation}
\max_\theta J(\theta)=\mathbb E_{\pi_\theta}
\left[\sum_{t=0}^{T-1}\gamma^t r_t\right].
\end{equation}
$T$ is the horizon and $\gamma$ the discount factor. Fit return targets and
bootstrap at nonterminal truncations. Reward forward tracking, uprightness,
and alternating support; penalize falls, slip, effort, and abrupt actions.
Revise rewards using physical diagnostics: contact events alone need not
represent useful steps.

Use an existing walker as a starting point for adaptation. Record whether a
candidate is pretrained, teacher-assisted, trained from random weights, or
resumed; these support different learning claims.

\section{Verify Motion Within a Fixed Budget}
Fix the task and budget before training. Freeze candidate weights and evaluate
continuous flat-ground trials from independently reset initial conditions.
Stop on a fall or the horizon without automatic resets. Inspect foot lifting,
load transfer, alternation, displacement, and slip alongside rendered motion.
Distinguish walking from pivoting, sliding, or repeated same-leg lifts.

Archive configurations, checkpoint hashes, training histories, and evaluations.
Stop at the criterion or budget; further training requires a separate decision.
Finite demonstrations do not establish general stability.

\section{Extend Toward Passive-Joint Motion}
Once walking is established, compare the unchanged reference with a movable
joint whose motor torque is disabled. Verify actual actuator effort. Adaptation
can then target compensation through remaining joints and contact dynamics.
Hardware transfer and broader robustness remain future work.

\begingroup
\raggedright
\bibliographystyle{IEEEtran}
\bibliography{references}
\endgroup
\end{document}
